% GENERATED FILE. Edit canonical lessons, then run npm run generate:books.
% canonical-source-hash: fnv1a64:58c76f52da02d6b9
% canonical-lessons: ZH-C206-lou, ZH-C206-ceng, ZH-C206-shanglou, ZH-C206-bangongshi, ZH-C206-qiantai

\chapter{Upstairs}
\label{ch:zh-upstairs}
\hlchaptermodality{206}

\begin{chapteropening}
\textbf{By the end of this chapter:} \emph{I can talk about buildings and floors, go upstairs, and find the office and the front desk.}

Complete the last lesson of chapter 206: I can talk about buildings and floors, go upstairs, and find the office and the front desk.
\end{chapteropening}

\section[lóu]{lóu — a building of several floors; a floor}

\label{lesson:ZH-C206-lou}



\begin{quote}
Before the new one: say the Mandarin for almost, then the Mandarin for distance.
\end{quote}

\subsection*{You'll want to know: lóu}
\emph{lóu} — \textquotedblleft{}a building of several floors; a floor\textquotedblright{}.

\begin{quote}
\emph{wǒ zhù de lóu hěn xīn} — The building I live in is new.
\end{quote}

Its characters come later; for now it is a word for the ear.

\begin{cousinweb}
\emph{lóu} is inherited from Middle Chinese \emph{luw} (Guangyun, in Baxter's transcription), glossed 'a house of more than one storey'.
\end{cousinweb}

\begin{grammarlens}[title={What you've built}]
One more word for asking where things are.
\end{grammarlens}

\subsection*{Your turn}
\begin{itemize}
\raggedright
  \item \emph{Say it:} \emph{lóu}
  \item \emph{Say it:} \emph{lóu}, once more
  \item \emph{Say it:} \emph{wǒ zhù de lóu hěn xīn} — The building I live in is new.
  \item \emph{Recall:} say the Mandarin for almost, then the Mandarin for distance, then say \emph{lóu} again
\end{itemize}

\begin{tcolorbox}[breakable,colback=teal!4,colframe=teal!35!black,title={Before you move on}]
What does lóu mean? (\textquotedblleft{}A building of several floors; a floor\textquotedblright{}.) Say it once more.
\end{tcolorbox}
\section[céng]{céng — floor; storey}

\label{lesson:ZH-C206-ceng}



\begin{quote}
Before the new one: say the Mandarin for distance, then the Mandarin for a building of several floors.
\end{quote}

\subsection*{You'll want to know: céng}
\emph{céng} — \textquotedblleft{}floor; storey\textquotedblright{}.

\begin{quote}
\emph{wǒmen de lóu yǒu èrshí céng} — Our building has twenty floors.
\end{quote}

Its characters come later; for now it is a word for the ear.

\begin{cousinweb}
\emph{céng} is inherited from Middle Chinese \emph{dzong} (Guangyun, in Baxter's transcription), glossed, like the form of \emph{lóu}, 'a house of more than one storey'.
\end{cousinweb}

\begin{grammarlens}[title={What you've built}]
One more word for asking where things are.
\end{grammarlens}

\subsection*{Your turn}
\begin{itemize}
\raggedright
  \item \emph{Say it:} \emph{céng}
  \item \emph{Say it:} \emph{céng}, once more
  \item \emph{Say it:} \emph{wǒmen de lóu yǒu èrshí céng} — Our building has twenty floors.
  \item \emph{Recall:} say the Mandarin for distance, then the Mandarin for a building of several floors, then say \emph{céng} again
\end{itemize}

\begin{tcolorbox}[breakable,colback=teal!4,colframe=teal!35!black,title={Before you move on}]
What does céng mean? (\textquotedblleft{}Floor; storey\textquotedblright{}.) Say it once more.
\end{tcolorbox}
\section[shànglóu]{shànglóu — to go upstairs}

\label{lesson:ZH-C206-shanglou}



\begin{quote}
Before the new one: say the Mandarin for a building of several floors, then the Mandarin for storey.
\end{quote}

\subsection*{You'll want to know: shànglóu}
\emph{shànglóu} — \textquotedblleft{}to go upstairs\textquotedblright{}.

\begin{quote}
\emph{wǒmen zuò diàntī shànglóu} — We're taking the lift upstairs.
\end{quote}

Its characters come later; for now it is a word for the ear.

\begin{cousinweb}
\emph{shànglóu} is built from \zh{上} 'to go up' and \emph{lóu} 'a building of several floors', both inherited from Middle Chinese, \emph{dzyangX} and \emph{luw} (Guangyun, in Baxter's transcription), which glosses the first 'to climb'.
\end{cousinweb}

\begin{grammarlens}[title={What you've built}]
One more word for asking where things are.
\end{grammarlens}

\subsection*{Your turn}
\begin{itemize}
\raggedright
  \item \emph{Say it:} \emph{shànglóu}
  \item \emph{Say it:} \emph{shànglóu}, once more
  \item \emph{Say it:} \emph{wǒmen zuò diàntī shànglóu} — We're taking the lift upstairs.
  \item \emph{Recall:} say the Mandarin for a building of several floors, then the Mandarin for storey, then say \emph{shànglóu} again
\end{itemize}

\begin{tcolorbox}[breakable,colback=teal!4,colframe=teal!35!black,title={Before you move on}]
What does shànglóu mean? (\textquotedblleft{}To go upstairs\textquotedblright{}.) Say it once more.
\end{tcolorbox}
\section[bàngōngshì]{bàngōngshì — office}

\label{lesson:ZH-C206-bangongshi}



\begin{quote}
Before the new one: say the Mandarin for storey, then the Mandarin for to go upstairs.
\end{quote}

\subsection*{You'll want to know: bàngōngshì}
\emph{bàngōngshì} — \textquotedblleft{}office\textquotedblright{}.

\begin{quote}
\emph{lǎoshī zài bàngōngshì} — The teacher is in the office.
\end{quote}

Its characters come later; for now it is a word for the ear.

\begin{cousinweb}
\emph{bàngōngshì} is built from \emph{bàn} 'to handle', \zh{公} 'public' and \emph{shì} 'a room', all inherited from Middle Chinese, \emph{beanH}, \emph{kuwng} and \emph{syit} (Guangyun, in Baxter's transcription), which glosses the first 'to provide': a room for public business.
\end{cousinweb}

\begin{grammarlens}[title={What you've built}]
One more word for asking where things are.
\end{grammarlens}

\subsection*{Your turn}
\begin{itemize}
\raggedright
  \item \emph{Say it:} \emph{bàngōngshì}
  \item \emph{Say it:} \emph{bàngōngshì}, once more
  \item \emph{Say it:} \emph{lǎoshī zài bàngōngshì} — The teacher is in the office.
  \item \emph{Recall:} say the Mandarin for storey, then the Mandarin for to go upstairs, then say \emph{bàngōngshì} again
\end{itemize}

\begin{tcolorbox}[breakable,colback=teal!4,colframe=teal!35!black,title={Before you move on}]
What does bàngōngshì mean? (\textquotedblleft{}Office\textquotedblright{}.) Say it once more.
\end{tcolorbox}
\section[qiántái]{qiántái — front desk; reception}

\label{lesson:ZH-C206-qiantai}



\begin{quote}
Before the new one: say the Mandarin for to go upstairs, then the Mandarin for office.
\end{quote}

\subsection*{You'll want to know: qiántái}
\emph{qiántái} — \textquotedblleft{}front desk; reception\textquotedblright{}.

\begin{quote}
\emph{yàoshi zài qiántái} — The key is at reception.
\end{quote}

Its characters come later; for now it is a word for the ear.

\begin{cousinweb}
\emph{qiántái} is built from \zh{前} 'front' and \emph{tái} 'a counter', both inherited from Middle Chinese, \emph{dzen} and \emph{doj} (Guangyun, in Baxter's transcription), which glosses the second as a raised platform.
\end{cousinweb}

\begin{grammarlens}[title={What you've built}]
One more word for asking where things are.
\end{grammarlens}

\subsection*{Your turn}
\begin{itemize}
\raggedright
  \item \emph{Say it:} \emph{qiántái}
  \item \emph{Say it:} \emph{qiántái}, once more
  \item \emph{Say it:} \emph{yàoshi zài qiántái} — The key is at reception.
  \item \emph{Recall:} say the Mandarin for to go upstairs, then the Mandarin for office, then say \emph{qiántái} again
\end{itemize}

\begin{tcolorbox}[breakable,colback=teal!4,colframe=teal!35!black,title={Before you move on}]
What does qiántái mean? (\textquotedblleft{}Front desk; reception\textquotedblright{}.) Say it once more.
\end{tcolorbox}
